\documentclass[11pt]{article}
\usepackage[margin=1in]{geometry}
\usepackage{amsmath,amssymb,amsthm}
\usepackage{algorithm}
\usepackage{algpseudocode}
\usepackage{hyperref}
\usepackage{booktabs}
\usepackage{xcolor}

\title{Recursive Phased-State Signatures:\\Problem Analysis and Mathematical Framework}
\author{Research note}
\date{May 28, 2026}

\newtheorem{definition}{Definition}
\newtheorem{lemma}{Lemma}
\newtheorem{theorem}{Theorem}
\newtheorem{problem}{Problem}

\begin{document}
\maketitle

\begin{abstract}
This note describes a research framework for a recursive phased-state
post-quantum signature scheme. The proposal is not a production-ready
cryptosystem. Its purpose is to turn a self-similar, fractal-inspired key
generation idea into a finite mathematical object with explicit security
games, implementation interfaces, and cryptanalytic questions. The conservative
instantiation is hash-based and uses recursive state expansion, phase-separated
commitments, one-time or few-time signature leaves, and optional batch
verification. More experimental finite-geometric variants can replace the
hash-only state expansion once the base framework is testable.
\end{abstract}

\section{Problem Analysis}

Current post-quantum signature options create several tradeoffs. ML-DSA is
practical but keeps the ecosystem heavily concentrated in structured lattices.
SLH-DSA is conservative and hash-based, but signatures are comparatively large.
Falcon/FN-DSA is compact, but implementation is more delicate. Additional NIST
signature candidates diversify assumptions, but several are complex,
assumption-heavy, or still exposed to active cryptanalysis.

The gap considered here is a signature framework with:
\begin{itemize}
  \item compact seed-derived private keys;
  \item explicitly separated signing phases;
  \item recursive, self-similar finite state generation;
  \item public commitments that bind all phases to one public key;
  \item a path to batch or aggregate verification;
  \item no dependence on floating point or real-valued chaotic maps.
\end{itemize}

\section{Solution Summary}

Let a root seed generate a tree of secret states. A phase corresponds to a depth
or labeled slice of that tree. Each state has a public projection. A public key
commits to all phase roots. A signature proves that one phase-local secret state
was authorized by the committed public key and signed a message.

\[
  s \longrightarrow S_\epsilon
    \longrightarrow \{S_v : v \in \mathcal{T}_{D,b}\}
    \longrightarrow \{P_v = \Pi_{\phi_{|v|}}(S_v)\}
    \longrightarrow PK.
\]

The first instantiation, RPS-H, is intentionally conservative: use SHAKE256 as a
domain-separated random oracle, a Merkle commitment per phase, and a one-time or
few-time signature layer at leaves. The recursive structure is evaluated for
compression, phase agility, and batch verification behavior before introducing
new hardness assumptions.

\section{Mathematical Framework}

\subsection{Parameters and Domains}

Let $\lambda \in \{128,192,256\}$ be the target security parameter. Let
$D \geq 1$ be the recursion depth and $b \geq 2$ be the branching factor. Define
the phase set
\[
  \Phi = \{\phi_0,\phi_1,\ldots,\phi_D\}.
\]

Let $\mathcal{T}_{D,b}$ be a rooted $b$-ary tree of depth $D$. A node is a path
$v=(j_1,\ldots,j_i)$ with $j_k \in \{0,\ldots,b-1\}$ and $|v|=i$. The root is
$\epsilon$.

Let
\[
  H : \{0,1\}^* \rightarrow \{0,1\}^{2\lambda}
\]
be modeled as a domain-separated random oracle or instantiated by SHAKE256.

\subsection{Recursive State Generation}

The root secret is sampled uniformly:
\[
  s \xleftarrow{\$} \{0,1\}^{\lambda}.
\]

The initial state is:
\[
  S_\epsilon = H(\textsf{RPS-root} \parallel s).
\]

For a child $v \parallel j$ at depth $i+1$, define:
\[
  S_{v \parallel j}
    = H(\textsf{RPS-child} \parallel S_v \parallel \phi_i \parallel j \parallel v).
\]

This rule is self-similar because the same expansion relation is used at every
scale. The phase label and path are included to prevent accidental related-key
reuse across phases and siblings.

\subsection{Public Projection and Phase Commitments}

For each node $v$ at phase $i=|v|$, define a public projection:
\[
  P_v = \Pi_{\phi_i}(S_v)
      = H(\textsf{RPS-public} \parallel \phi_i \parallel v \parallel S_v).
\]

Let $L_i=\{P_v : |v|=i\}$ be the set of public projections in phase $i$. Define
the phase commitment:
\[
  C_i = \mathsf{MerkleRoot}(L_i).
\]

The global public key is:
\[
  PK = H(\textsf{RPS-pk} \parallel C_0 \parallel C_1 \parallel \cdots \parallel C_D).
\]

The private key is either the seed $s$ or the seed plus cached internal states:
\[
  SK = s
  \quad\text{or}\quad
  SK = (s,\{S_v\}_{v \in \mathcal{C}}).
\]

\section{Algorithms}

\begin{algorithm}
\caption{$\mathsf{KeyGen}(1^\lambda,D,b)$}
\begin{algorithmic}[1]
\State $s \xleftarrow{\$} \{0,1\}^{\lambda}$
\State $S_\epsilon \gets H(\textsf{RPS-root} \parallel s)$
\For{$i=0$ to $D-1$}
  \ForAll{$v$ such that $|v|=i$}
    \For{$j=0$ to $b-1$}
      \State $S_{v\parallel j} \gets H(\textsf{RPS-child}\parallel S_v\parallel\phi_i\parallel j\parallel v)$
    \EndFor
  \EndFor
\EndFor
\For{$i=0$ to $D$}
  \State $L_i \gets \{H(\textsf{RPS-public}\parallel\phi_i\parallel v\parallel S_v): |v|=i\}$
  \State $C_i \gets \mathsf{MerkleRoot}(L_i)$
\EndFor
\State $PK \gets H(\textsf{RPS-pk}\parallel C_0\parallel\cdots\parallel C_D)$
\State \Return $(PK,SK=s)$
\end{algorithmic}
\end{algorithm}

\begin{algorithm}
\caption{$\mathsf{Sign}(SK,m,i)$}
\begin{algorithmic}[1]
\State Select an unused path $v$ with $|v|=i$
\State Derive $S_v$ from $s$ along path $v$
\State $x_v \gets H(\textsf{RPS-ots-secret}\parallel S_v\parallel H(m))$
\State $\sigma_{\mathsf{ots}} \gets \mathsf{OTS.Sign}(x_v,m)$
\State $P_v \gets H(\textsf{RPS-public}\parallel\phi_i\parallel v\parallel S_v)$
\State $\pi_v \gets \mathsf{MerkleAuthPath}(P_v,C_i)$
\State \Return $\sigma=(i,v,\sigma_{\mathsf{ots}},P_v,\pi_v)$
\end{algorithmic}
\end{algorithm}

\begin{algorithm}
\caption{$\mathsf{Verify}(PK,m,\sigma)$}
\begin{algorithmic}[1]
\State Parse $\sigma=(i,v,\sigma_{\mathsf{ots}},P_v,\pi_v)$
\If{$|v|\neq i$} \State \Return $\mathsf{reject}$ \EndIf
\If{$\mathsf{OTS.Verify}(P_v,m,\sigma_{\mathsf{ots}})=\mathsf{false}$}
  \State \Return $\mathsf{reject}$
\EndIf
\If{$\mathsf{MerkleVerify}(P_v,\pi_v,C_i)=\mathsf{false}$}
  \State \Return $\mathsf{reject}$
\EndIf
\If{$PK \neq H(\textsf{RPS-pk}\parallel C_0\parallel\cdots\parallel C_D)$}
  \State \Return $\mathsf{reject}$
\EndIf
\State \Return $\mathsf{accept}$
\end{algorithmic}
\end{algorithm}

\section{Security Definitions}

\begin{definition}[Cross-Scale One-Wayness]
Given public projections
\[
  \mathcal{P}=\{P_v : v \in Q \subseteq \mathcal{T}_{D,b}\},
\]
an adversary should not recover any corresponding secret state $S_u$ for
$u\notin Q_{\mathsf{revealed}}$ with probability non-negligibly better than
$2^{-\lambda}$, even when projections from parents, children, and siblings are
available.
\end{definition}

\begin{definition}[Phase Binding]
A signature created under phase $\phi_i$ must not verify under phase $\phi_j$
for $i\neq j$. Equivalently, the verification relation must bind
\[
  (PK,m,\sigma,i,v)
\]
to exactly one phase label and path.
\end{definition}

\begin{definition}[RPS-EUF-CMA]
The unforgeability game gives the adversary $PK$ and signing access for chosen
messages and phases, subject to path-use rules. The adversary wins if it outputs
$(m^\star,\sigma^\star)$ accepted by $\mathsf{Verify}$ such that the signed
message/path pair was not previously returned by the signing oracle.
\end{definition}

\section{Proof Sketch for the Conservative Hash-Based Variant}

\begin{lemma}[Merkle Membership Binding]
If the Merkle tree hash is collision resistant, then an adversary cannot replace
a committed projection $P_v$ with a distinct $P'_v$ under the same phase root
$C_i$ except with negligible probability.
\end{lemma}

\begin{proof}
Any successful replacement induces either an equal leaf hash for distinct leaves
or a collision at some internal node along the authentication path. Both cases
contradict collision resistance of the Merkle compression/hash function.
\end{proof}

\begin{lemma}[Phase Separation]
If $H$ is domain separated and modeled as a random oracle, then values derived
under distinct labels $\phi_i\neq\phi_j$ are independent oracle domains.
\end{lemma}

\begin{proof}
All calls include explicit strings and phase labels. Therefore an oracle query
for one phase does not answer a query in another phase unless the full input
strings collide. With canonical encoding, this requires identical labels and
paths.
\end{proof}

\begin{theorem}[RPS-H Unforgeability, Informal]
Assume the one-time signature is EUF-CMA secure for unused leaves, the Merkle
commitment is collision resistant, and $H$ behaves as a domain-separated random
oracle. Then RPS-H is EUF-CMA secure up to the probability of OTS forgery,
Merkle collision, or reuse of a previously exposed signing state.
\end{theorem}

\begin{proof}
Let an adversary output a valid forgery. If its leaf projection is not committed
under the claimed phase root, Merkle verification should fail, except by the
membership-binding lemma. If the projection is committed but the leaf was not
previously used, then the adversary has forged the OTS. If the forgery changes
the phase, the phase-separation lemma prevents reinterpreting the same proof
under another phase. If the adversary reuses an exposed leaf, the attack is
outside the one-time use condition and must be prevented by state management or
by replacing OTS with a stateless many-leaf construction.
\end{proof}

\section{Finite-Geometric Experimental Variant}

After RPS-H is implemented, the recursive expansion can be generalized to a
finite geometric state space. Let $S_v \in \mathbb{F}_q^n$. Define:
\[
  S_{v\parallel j}
  = H_{\mathbb{F}_q^n}
    \left(A_{\phi_i,j}S_v + B_{\phi_i,j}
      + N_{\phi_i,j}(S_v) \bmod q\right),
\]
where $A_{\phi_i,j}\in\mathbb{F}_q^{n\times n}$, $B_{\phi_i,j}\in\mathbb{F}_q^n$,
and $N_{\phi_i,j}$ is a low-degree nonlinear map.

\begin{problem}[Hidden Recursive Embedding Problem]
Given public projections $\{P_v\}$ of a recursively generated finite geometric
object, recover a secret state $S_u$, a valid hidden path, or a trapdoor relation
linking two phases.
\end{problem}

This variant is only credible if it survives attacks based on linearization,
Gröbner bases, cycle finding, meet-in-the-middle search, related-key analysis,
and quantum amplitude amplification.

\section{Implementation Framework}

\begin{center}
\begin{tabular}{lll}
\toprule
Component & Conservative choice & Experimental replacement \\
\midrule
Expansion & SHAKE256 XOF & finite-field recursive map \\
Projection & hash of state & lattice/code/geometric projection \\
Commitment & Merkle tree & vector commitment or accumulator \\
Leaf signing & WOTS+/FORS-like OTS & proof of hidden relation \\
State mode & stateful counters & stateless hypertree sampling \\
Batching & shared Merkle/OTS checks & aggregate verification relation \\
\bottomrule
\end{tabular}
\end{center}

The minimum implementation API is:
\[
\begin{aligned}
  &(PK,SK) \leftarrow \mathsf{RPS.KeyGen}(\lambda,D,b),\\
  &\sigma \leftarrow \mathsf{RPS.Sign}(SK,m,i),\\
  &\{0,1\} \leftarrow \mathsf{RPS.Verify}(PK,m,\sigma),\\
  &\{0,1\} \leftarrow \mathsf{RPS.BatchVerify}(PK,\{m_k,\sigma_k\}_k).
\end{aligned}
\]

\section{Evaluation Plan}

The first prototype should measure:
\begin{itemize}
  \item public key size as a function of $D$ and $b$;
  \item private key size with seed-only and cached modes;
  \item signature size by phase;
  \item signing and verification time;
  \item batch verification savings;
  \item failure behavior under phase misuse and leaf reuse;
  \item cross-phase statistical and structural leakage.
\end{itemize}

\section{Conclusion}

Recursive phased-state signatures are best treated as a framework first and a
new hardness assumption second. The conservative hash-based version makes the
architecture concrete while relying on familiar assumptions. The more ambitious
finite-geometric version should be introduced only after the recursive framework
has clear security games, benchmarks, and known attack targets.

\section*{References}

\begin{itemize}
  \item NIST, Post-Quantum Cryptography Standardization Project:
  \url{https://csrc.nist.gov/Projects/Post-Quantum-Cryptography}
  \item NIST, Round 3 Additional Digital Signature Schemes:
  \url{https://csrc.nist.gov/projects/pqc-dig-sig/round-3-additional-signatures}
  \item NIST IR 8610, Status Report on the Second Round of the Additional
  Digital Signature Schemes for the NIST Post-Quantum Cryptography
  Standardization Process, 2026.
\end{itemize}

\end{document}
